% JETSPIN jet random noise documentation

\section{Random noise\label{sec:Random-noise}}

In JETSPIN, it is possible to add a random noise to Eq. \ref{eq:force-EOM}
in order to model the instability at the nozzle triggering the bending
motion. This procedure can be used as an alternative to (or in junction
with) the strategy presented in Section \ref{sec:Perturbations-at-the}.

In particular, collected in $\vec{\textbf{f}}_{old,i}=\vec{\textbf{f}}_{el,i}+\vec{\textbf{f}}_{c,i}+\vec{\textbf{f}}_{\upsilon e,i}+\vec{\textbf{f}}_{st,i}+\vec{\textbf{f}}_{g,i}+\vec{\textbf{f}}_{diss,i}+\vec{\textbf{f}}_{rand}$
all the force terms, we modify Eq. \ref{eq:force-EOM} as
\begin{equation}
m_{i}\frac{d\vec{\boldsymbol{\upsilon}}_{i}}{dt}=\vec{\textbf{f}}_{old,i}-m_{i}\gamma_{rand}\vec{\boldsymbol{\upsilon}}_{i}+\sqrt{2m_{i}^{2}D_{rand}}\vec{\boldsymbol{\eta}_{i}}(t),
\end{equation}

following a standard Langevin approach, where $D_{rand}$ is a diffusion
coefficient in the velocity space (in cgs units, $\text{cm}^{2}/\text{s}^{3}$)
and as $\gamma_{rand}$ a dissipative term ($\text{s}^{-1}$). We
highlight that $D_{rand}$ sets the width of the fluctuations. In
particular, it is possible to demonstrate \cite{gillespie2012simple}
that the variance $\sigma_{rand}\left(t\right)^{2}$ of the velocity
$\vec{\boldsymbol{\upsilon}}_{i}$ ($\text{cm}^{2}/\text{s}^{2}$)
due to only the random and dissipative terms at time $t$, computed
over an ensemble of stochastic trajectories, is equal to $\sigma_{rand}\left(t\right)^{2}\equiv$$<\left[\vec{\boldsymbol{\upsilon}}_{i}-<\vec{\boldsymbol{\upsilon}}_{i}>\right]^{2}>$
$=\left(D_{rand}/\gamma_{rand}\right)\left(1-e^{-2\gamma_{rand}t}\right)$,
and, consequently, $\lim\limits _{t\rightarrow\infty}\sigma_{rand}\left(t\right)^{2}=D_{rand}/\gamma_{rand}$,
which is sometimes called fluctuation-dissipation Langevin relation.
In \texttt{input} file, the user should specified both $D_{rand}$
and $\sigma_{rand}{}^{2}$. The $\gamma_{rand}$ is automatically
set to be equal to $\gamma_{rand}=D_{rand}/\sigma_{rand}{}^{2}$ by
exploiting the Langevin relation.
